\documentclass{training}
\schoollogo{../../resources/intros/UI-03-building-an-interface.jpg}
\setlang{french}
\newcommand{\uxfond}{../../resources/fond/UI-03-building-an-interface.png}
% Filigrane PAR MODULE : \uxfond (PNG A4 dilue 30 %). Cadre noir vectoriel commun.
% Proportions original 744 px : marge 20/744, trait 1/744. Exterieur BLANC.
% Cadre via eso-pic +\put (coordonnees absolues) — pas de remember picture (decalage).
\usepackage{eso-pic}

\providecommand{\uxfond}{../../resources/fond/norme_fond_paliers_30.png}

\newlength{\uxmargin}
\newlength{\uxthick}
\newlength{\uxinnerw}
\newlength{\uxinnerh}
\setlength{\uxmargin}{0.02688\paperwidth}% 20/744
\setlength{\uxthick}{0.001344\paperwidth}% 1/744
\setlength{\uxinnerw}{\dimexpr\paperwidth-2\uxmargin-2\uxthick\relax}
\setlength{\uxinnerh}{\dimexpr\paperheight-2\uxmargin-2\uxthick\relax}

% Contenu strictement a l'interieur du trait (+ jour confortable).
\geometry{%
  includeheadfoot,
  left=22mm,
  right=22mm,
  top=20mm,
  bottom=18mm,
  headheight=30pt,
  headsep=5mm,
  footskip=12mm%
}

\AddToShipoutPictureBG{%
  \AtPageLowerLeft{%
    % Fond blanc page entiere
    \put(0,0){\color{white}\rule{\paperwidth}{\paperheight}}%
    % Filigrane module a l'interieur du cadre
    \put(\LenToUnit{\dimexpr\uxmargin+\uxthick\relax},\LenToUnit{\dimexpr\uxmargin+\uxthick\relax}){%
      \includegraphics[width=\uxinnerw,height=\uxinnerh]{\uxfond}}%
    % Cadre noir : bas, haut, gauche, droite
    \put(\LenToUnit{\uxmargin},\LenToUnit{\uxmargin}){%
      \color{black}\rule{\dimexpr\paperwidth-2\uxmargin\relax}{\uxthick}}%
    \put(\LenToUnit{\uxmargin},\LenToUnit{\dimexpr\paperheight-\uxmargin-\uxthick\relax}){%
      \color{black}\rule{\dimexpr\paperwidth-2\uxmargin\relax}{\uxthick}}%
    \put(\LenToUnit{\uxmargin},\LenToUnit{\uxmargin}){%
      \color{black}\rule{\uxthick}{\dimexpr\paperheight-2\uxmargin\relax}}%
    \put(\LenToUnit{\dimexpr\paperwidth-\uxmargin-\uxthick\relax},\LenToUnit{\uxmargin}){%
      \color{black}\rule{\uxthick}{\dimexpr\paperheight-2\uxmargin\relax}}%
  }%
}

\makeatletter
\fancyhf{}
\fancyheadoffset{0pt}
\renewcommand{\headrulewidth}{0pt}
\lhead{\parbox[t]{0.56\textwidth}{\RaggedRight\footnotesize\sloppy\@title}}
\rhead{\parbox[t]{0.36\textwidth}{\RaggedLeft\footnotesize\sloppy\@subtitle}}
\cfoot{\thepage}
% Coupures : underscore et slash dans les noms de fichiers / consignes
\DeclareRobustCommand{\_}{\textunderscore\allowbreak}
\makeatother

% Evite les depassements de longues lignes (Github, Name_FirstName_...)
\sloppy
\emergencystretch=3em
\hbadness=10000

% Sommaire type 42 : chapitres seulement
\setcounter{tocdepth}{0}
\renewcommand{\cftchappresnum}{}
\renewcommand{\cftchapaftersnum}{}
\setlength{\cftchapnumwidth}{2.0em}
\setlength{\cftbeforechapskip}{0.35em}
\cftsetindents{chapter}{0em}{2.0em}


\begin{document}

\title{Construire une interface}
\subtitle{Piscine UX/UI --- UI}
\member{Piscine UX/UI}{https://github.com/01-edu/UX-UI}
\summary{Un styleguide, une bibliothèque, des écrans, du mouvement, un test de désirabilité.}
\version{1.0}
\prerequisite{Le module précédent de cette piscine, ou aucun si c'est le premier.}
\maketitle

\newpage
\tableofcontents
\newpage


\chapter{Instructions}

L'avant-propos n'est pas un exercice. Il ne sera pas noté. Les exercices commencent au chapitre suivant, du contrat le plus simple au plus dur. Un exercice plus tardif n'est pas pris en compte si un plus tôt ne tient pas.

On ne regarde que ce qui est dans le répertoire de rendu de l'exercice. N'ajoutez pas un fichier que l'exercice ne nomme pas. Le nom d'un livrable suit le motif imprimé dans l'exercice : \texttt{Name\_FirstName\_DeliverableName\_Date\_VersionNumber}.

Le geste ne vous est pas donné. Les exemples d'un exercice portent parfois une contrainte que la phrase ne répète pas. Lisez-les. Vos autres références sont le manuel de l'outil choisi, la personne assise à côté, et les liens imprimés dans l'exercice.

L'ordre est celui de la quête. Sauter une étape parce qu'un écran plus tard semble plus intéressant produit un écran que personne ne sait expliquer.

\subsection*{Des images}

La marque en couverture est l'image de ce module. Les paliers pâles derrière le texte sont le fond de page. Ni l'un ni l'autre n'est un indice.

\subsection*{Des machines qui écrivent}

Vous pouvez vous servir d'un outil pour vérifier un mot, comme d'un dictionnaire. Vous ne pouvez pas remplacer un entretien, un test ou un fichier par un texte généré et le signer. Si vous ne pouvez pas montrer la personne, l'écran ou la ligne que vous avancez, l'exercice n'est pas fait.


\chapter{Avant-propos}

\emph{Objectified} de Gary Hustwit (2009) filme des gens qui se disputent sur une poignée, un cure-dent, une chaise. La dispute est le travail. Cette quête est cette dispute, tenue sur des écrans que vous dessinez. Une seconde étagère est dans l'avant-propos, une playlist sur l'IA. Elle est là pour entendre comment parlent les outils. Elle n'autorise pas à rendre un deck généré.

\begin{center}
\includegraphics[width=0.78\textwidth,height=0.32\textheight,keepaspectratio]{../../resources/intros/UI-03-building-an-interface.jpg}\\[0.4em]
{\footnotesize Image : freddie marriage fredmarriage, CC0}
\end{center}

À regarder. Les vidéos sont un contexte. Elles ne sont pas le livrable.
\begin{itemize}
\item UI - Introduction \& Quest n°1, 01Edu \url{https://www.youtube.com/watch?v=UbE6WMeuvQg&list=PLHyAJ_GrRtf-_aBft2jQL4-QAunn2bLvu}
\item Claude AI Enablement piscine, 01Edu \url{https://www.youtube.com/watch?v=M94thLMc6vA&list=PLCiJelRloNB8}
\end{itemize}
La dernière playlist est une piscine d'IA (\emph{Claude AI Enablement}). Regardez-la si vous voulez entendre comment on vend un modèle comme assistant de design. Les fichiers rendus viennent encore de votre recherche, de vos écrans et de vos tests. Une planche générée que vous ne savez pas raconter ne compte pas.

\chapter{Exercice 00 : UI IV - Ex 1 - Styleguide}
\exnumber{00}
\extitle{UI IV - Ex 1 - Styleguide}
\exfiles{ex00/fichier Figma, nomme Name\_FirstName\_DeliverableName\_Date\_VersionNumber}
\exauthorize{Those named in the exercise}
\makeheaderfiles

\textbf{Contexte :}

In this quest, you'll have to build your own interface.

After designing your own style guide with a moodboard, colors and typography, you can use the series of consistency principles, atomic design structure, the techniques to build you own library and then a series of screens. Step by step, you'll have built an animated high-fidelity prototype.

The very last step will be to test it on users to detect the impressions they get when they discover your app.

Once again, feel free to get as much inspiration as you need from existing design systems. Dribbble, consistency and structure are your best friends!

\textbf{Instructions :}

Your agency works on a digital product to create an e-commerce app for second-hand electronic devices. You collaborate with a UX designer from your team. They have deisgned the wireframes and pass it on to you!

Design a style guide (= a moodboard with typography + examples of typography with buttons, title, description and content).

Also think about the values and spirit you want to spread through this moodboard, we'll get back to it in Ex. 5.

\textbf{Livrables :}

\begin{itemize}
\item A Figma moodboard that includes:
\begin{itemize}
\item Des photos
\item 2 to 3 main colors
\item Examples of typography applied to graphic elements
\begin{itemize}
\item Title
\item Description or subtitle
\item Content
\item Button
\end{itemize}
\end{itemize}
\item Déposez le fichier Figma sur Github, nommé ainsi : « Name\_FirstName\_DeliverableName\_Date\_VersionNumber ». Exemple : « Doe\_John\_MidFiPrototype\_05242024\_V1 ».
\item Ajoutez un titre dans le fichier.
\item Soignez l'aspect global du fichier. Il doit être clair, simple et lisible. Vous pouvez regarder canva.com, sans surcharger le dessin de détails.
\end{itemize}

\textbf{Ressources :}

\begin{itemize}
\item Video tutorial on Figma \url{https://www.youtube.com/watch?v=FTFaQWZBqQ8}
\item Beginning graphic design \url{https://edu.gcfglobal.org/en/beginning-graphic-design/}
\item Gareth graphic studio \url{https://www.youtube.com/c/GarethDavidStudio}
\item Example of a Style Guide in UI \url{https://www.pinterest.pt/pin/454019206175050454/}
\end{itemize}


\chapter{Exercice 01 : UI IV - Ex 2 - Bibliothèque de design system}
\exnumber{01}
\extitle{UI IV - Ex 2 - Bibliothèque de design system}
\exfiles{ex01/fichier Figma, nomme Name\_FirstName\_DeliverableName\_Date\_VersionNumber}
\exauthorize{Ceux nommés dans l'exercice}
\makeheaderfiles

\textbf{Instructions :}

À partir du moodboard conçu à l'Ex 1, créez un design system avec au moins 30 éléments, incluant atoms, molecules et organisms, et incluez des multistates.

\textbf{Livrables :}

\begin{itemize}
\item Un fichier Figma avec une bibliothèque comprenant :
\begin{itemize}
\item Au moins 30 éléments de design graphique
\item Au moins 5 atomes
\item Au moins 5 molécules
\item Au moins 5 organismes
\item Au moins 10 variations multi-états
\end{itemize}
\item Déposez le fichier Figma sur Github, nommé ainsi : « Name\_FirstName\_DeliverableName\_Date\_VersionNumber ». Exemple : « Doe\_John\_MidFiPrototype\_05242024\_V1 ».
\item Ajoutez un titre dans le fichier.
\item Soignez l'aspect global du fichier. Il doit être clair, simple et lisible. Vous pouvez regarder canva.com, sans surcharger le dessin de détails.
\end{itemize}

\textbf{Pistes :}

\begin{itemize}
\item Commencez par vous demander quels éléments clés peuvent apparaître dans une application e-commerce.
\item Inspirez-vous des design systems existants. Parcourez Dribbble et la communauté Figma.
\end{itemize}

\textbf{Ressources :}

\begin{itemize}
\item FICHIER FIGMA (à dupliquer) \url{https://www.figma.com/file/U26mctkrOqxiUeAlNHzSxW/UI-IV---Ex-3---E-COMMERCE%3A-Design-screens!?node-id=0%3A1}
\item Introduction aux design systems \url{https://www.youtube.com/watch?v=xEuBCUngJ_U}
\item Créer un design system - cours Figma \url{https://www.youtube.com/watch?v=RYDiDpW2VkM}
\item Construire une library de design system - Lyft \url{https://medium.com/tap-to-dismiss/building-a-design-system-library-532ef2492811}
\item Exemples de library de boutons UI \url{https://www.pinterest.co.kr/pin/663014376380850291/}
\end{itemize}


\chapter{Exercice 02 : UI IV - Ex 3 - Écrans}
\exnumber{02}
\extitle{UI IV - Ex 3 - Écrans}
\exfiles{ex02/fichier Figma, nomme Name\_FirstName\_DeliverableName\_Date\_VersionNumber}
\exauthorize{Ceux nommés dans l'exercice}
\makeheaderfiles

\textbf{Instructions :}

À partir des wireframes et de votre library de design system, amenez votre prototype moyenne fidélité à un niveau haute fidélité !

Sélectionnez 5 écrans qui montreront le user flow d'un client qui :

\begin{itemize}
\item arrive sur la page d'accueil
\item parcourt le contenu
\item consulte une page produit
\item ajoute le produit au panier
\item finalise l'achat
\end{itemize}

Utilisez vos composants, call to actions et visuels pour compléter les wireframes.

\textbf{Livrables :}

\begin{itemize}
\item 5 écrans haute fidélité depuis Figma
\item Déposez le fichier Figma sur Github, nommé ainsi : « Name\_FirstName\_DeliverableName\_Date\_VersionNumber ». Exemple : « Doe\_John\_MidFiPrototype\_05242024\_V1 ».
\item Ajoutez un titre dans le fichier.
\item Soignez l'aspect global du fichier. Il doit être clair, simple et lisible. Vous pouvez regarder canva.com, sans surcharger le dessin de détails.
\end{itemize}

\textbf{Ressources :}

\begin{itemize}
\item Unsplash \url{https://unsplash.com/} pour les visuels
\item De la basse fidélité à la haute fidélité \url{https://mockitt.wondershare.com/prototyping/low-fidelity-vs-high-fidelity.html}
\end{itemize}


\chapter{Exercice 03 : UI IV - Ex 4 - Animation}
\exnumber{03}
\extitle{UI IV - Ex 4 - Animation}
\exfiles{ex03/fichier Figma, nomme Name\_FirstName\_DeliverableName\_Date\_VersionNumber}
\exauthorize{Ceux nommés dans l'exercice}
\makeheaderfiles

\textbf{Instructions :}

Animez vos 5 écrans sur Figma. Pensez au multistate !!

\textbf{Livrable :}

\begin{itemize}
\item Un prototype haute fidélité animé sur Figma
\item Déposez le fichier Figma sur Github, nommé ainsi : « Name\_FirstName\_DeliverableName\_Date\_VersionNumber ». Exemple : « Doe\_John\_MidFiPrototype\_05242024\_V1 ».
\item Ajoutez un titre dans le fichier.
\item Soignez l'aspect global du fichier. Il doit être clair, simple et lisible. Vous pouvez regarder canva.com, sans surcharger le dessin de détails.
\end{itemize}

\textbf{Ressources :}

\begin{itemize}
\item Gérer efficacement vos designs - Comment Mockplus State Page aide à voir et gérer plusieurs états \url{https://help.mockplus.com/p/372}
\end{itemize}

\textbf{Citation :}

« Parlez humain : utilisez les mots que vos utilisateurs utilisent. » Dan Saffer | Microinteractions


\chapter{Exercice 04 : UI IV - Ex 5 - Test de désirabilité}
\exnumber{04}
\extitle{UI IV - Ex 5 - Test de désirabilité}
\exfiles{ex04/archive zip, nomme Name\_FirstName\_DeliverableName\_Date\_VersionNumber}
\exauthorize{Those named in the exercise}
\makeheaderfiles

\textbf{Instructions :}

Run a desirability test!

The point of this phase is to make sure the values and the spirit you chose for your design.

\begin{itemize}
\item Get back to the values and spirit you chose in Ex 1. Extract 2 to 5 adjectives.
\item Prepare a board with 25+ adjectives, randomly distributed.
\item Show your design to 5 different people, and ask them to select the adjectives among the list from the board that they are inspired with.
\item If the selected adjectives chosen by the users are consistent with the ones you had chosen, congratulations!
\item If the selected adjectives chosen by the users are not consistent, that's fine! Make 3 suggestions about how to make your moodboard and library evolve.
\end{itemize}

\textbf{Livrables :}

\begin{itemize}
\item A desirability testing board with 25+ different adjectives
\item Pictures from the desirability test
\item A PDF document with the outcomes of the desirability test
\item If the outcomes are not consistent with the values, 3 suggestions on how to make the moodboard and design library evolve.
\item Déposez les documents sur Github, dans un zip nommé « ProjectTitle\_Name\_FirstName », et chaque livrable ainsi : « Name\_FirstName\_DeliverableName\_Date\_VersionNumber ». Exemple : « Doe\_John\_MidFiPrototype\_05242024\_V1 ».
\item Ajoutez un titre dans tous les documents écrits.
\item Soignez l'aspect global des livrables. Ils doivent être clairs, simples et lisibles. Vous pouvez regarder canva.com, sans surcharger le dessin de détails.
\end{itemize}

\textbf{Pistes :}

\begin{itemize}
\item Avoid confirmation biais: don't tell them the values and ask them to approve them!
\item Good to know: Desirability testing is great to assess the impressions a design gives to the users. Whereas Usability testing is great to understand what users understand your prototype and how they navigate across the product.
\end{itemize}

\textbf{Ressources :}

\begin{itemize}
\item Desirability Testing analyzing an emotional response to a design \url{https://www.slideshare.net/megangrocki/desirability-testing-analyzing-emotional-response-to-a-design-11262575}
\end{itemize}

\textbf{Citation :}

\begin{itemize}
\item 'Interaction design focuses on the design of behavior of a person's interaction with a digital interface. It is also concerned with satisfying the needs and desires of the people who will interact with a product or service.' Alan Cooper, About Face: The Essentials of Interaction Design
\end{itemize}


\end{document}
